% SPDX-License-Identifier: CC-BY-4.0
\section{Proof of the main results}\label{sec:main}

\begin{proof}[Proof of Theorem~\ref{thm:main}]
Theorem~\ref{thm:construction} gives a polynomial with the required
multiplicities and degree at most \(\Phi(n,k,\ell)\), and
Theorem~\ref{thm:per-order-lower} shows that every such polynomial has degree at
least \(\Phi(n,k,\ell)\). So the constructed polynomial has degree exactly
\(\Phi(n,k,\ell)\).
\end{proof}

\begin{lemma}[The minimum over the origin order]\label{lem:minimum}
For \(n\ge1\) and \(k\ge1\), \(\min_{0\le\ell<k}\Phi(n,k,\ell)\) equals
\(2k+n-2-\lfloor\log_2k\rfloor\) for \(k<2^{n-1}\), and
\(2k-\lfloor k/2^{n-1}\rfloor\) for \(k\ge2^{n-1}\)
\textnormal{\small\leanref{claims/PerOrderMinimumDirect.lean}{[Lean: this proof]}\,\leanref{claims/PerOrderMinimum.lean}{[Lean: via Appendix~\ref{sec:second-proof}]}}.
\end{lemma}

\begin{proof}
Write \(m=k-\ell-1=q2^n+r\) as before and \(h(m)=2q+s_2(r)\). By
Lemma~\ref{lem:legendre}, \(\Phi=n+2k-2-h(m)\), so the minimum is
\(n+2k-2-\max_{0\le m\le k-1}h(m)\). Write \(k-1=Q2^n+R\) with \(0\le R<2^n\).
Let \(\lambda=\lfloor\log_2(R+1)\rfloor\). The largest digit sum among
\(0,\ldots,R\) is \(\lambda\): it is attained at \(2^\lambda-1\le R\), and every
number up to \(R\) is below \(2^{\lambda+1}-1\), so it has at most \(\lambda\)
binary ones. With \(q=Q\) the best value of \(h\) is \(2Q+\lambda\). With
\(q\le Q-1\), possible only if \(Q\ge1\), it is at most \(2Q-2+n\), attained at
\(m=Q2^n-1\).
\begin{itemize}
 \item If \(Q=0\), that is \(k\le2^n\), the minimum is
 \(n+2k-2-\lfloor\log_2k\rfloor\). This is the first case when \(k<2^{n-1}\). For
 \(2^{n-1}\le k<2^n\) it is \(2k-1\), and for \(k=2^n\) it is \(2k-2\); both equal
 \(2k-\lfloor k/2^{n-1}\rfloor\).
 \item If \(Q\ge1\), then \(k=Q2^n+(R+1)\) with \(1\le R+1\le2^n\), and
 \(\lfloor k/2^{n-1}\rfloor=2Q+\lfloor(R+1)/2^{n-1}\rfloor\). Moreover
 \(\max(\lambda,n-2)=n-2+\lfloor(R+1)/2^{n-1}\rfloor\) in each of
 the three ranges \(R+1<2^{n-1}\), \(2^{n-1}\le R+1<2^n\) and \(R+1=2^n\) (for
 \(n=1\) the first is empty). So the minimum is
 \(n+2k-2-2Q-(n-2)-\lfloor(R+1)/2^{n-1}\rfloor=2k-\lfloor k/2^{n-1}\rfloor\). \qedhere
\end{itemize}
\end{proof}

\begin{proof}[Proof of Corollary~\ref{cor:main}]
Since \(D(n,k)=\min_\ell\delta(n,k,\ell)\), Theorem~\ref{thm:main} and
Lemma~\ref{lem:minimum} give the value. By Proposition~\ref{prop:cover-bound}, or
directly from the value, every \(k\)-admissible polynomial has degree at least
\(2k-\lfloor k/2^{n-1}\rfloor\). So a polynomial of exactly this degree exists if
and only if \(D(n,k)\) equals it. Compare the two in three cases.
\begin{itemize}
 \item If \(k\ge2^{n-1}\), they are equal.
 \item If \(2^{n-2}\le k<2^{n-1}\), then \(\lfloor\log_2k\rfloor=n-2\) and
 \(\lfloor k/2^{n-1}\rfloor=0\), so both equal \(2k\).
 \item If \(k<2^{n-2}\), then \(n\ge3\), \(\lfloor\log_2k\rfloor\le n-3\) and
 \(\lfloor k/2^{n-1}\rfloor=0\), so \(D(n,k)\ge2k+1\). \qedhere
\end{itemize}
\end{proof}

\begin{corollary}\label{cor:extended}
For \(n\ge1\) and \(1\le k<3\cdot2^{n-1}\), \(D(n,k)=2k+n-2-\lfloor\log_2k\rfloor\)
\lean{claims/BinaryMultiplicityDegreeExtendedRange.lean}.
\end{corollary}

\begin{proof}
For \(k<2^{n-1}\) this is Corollary~\ref{cor:main}. Otherwise compare with
\(2k-\lfloor k/2^{n-1}\rfloor\).
\begin{itemize}
 \item If \(2^{n-1}\le k<2^n\), then \(\lfloor\log_2k\rfloor=n-1\) and
 \(\lfloor k/2^{n-1}\rfloor=1\), so both expressions equal \(2k-1\).
 \item If \(2^n\le k<3\cdot2^{n-1}\), then \(\lfloor\log_2k\rfloor=n\) and
 \(\lfloor k/2^{n-1}\rfloor=2\), so both equal \(2k-2\). \qedhere
\end{itemize}
\end{proof}

At \(k=3\cdot2^{n-1}\) the logarithmic expression is \(2k-2\), while
\(D(n,k)=2k-3\).

\begin{remark}[Which origin orders are extremal]\label{rem:extremal}
The orders \(\ell\) attaining \(D(n,k)\) are those for which \(k-\ell-1=q2^n+r\)
maximizes \(2q+s_2(r)\). For \(k\le2^n\) they are the \(\ell\) with
\(s_2(k-\ell-1)=\lfloor\log_2k\rfloor\), for instance
\(\ell=k-2^{\lfloor\log_2k\rfloor}\) \lean{claims/PerOrderMinimumDirect.lean}.
\end{remark}

\subsection*{Fields of characteristic two}

The values of Theorem~\ref{thm:main} and Corollary~\ref{cor:main} hold over every
field of characteristic two, by a descent to \(\Ft\). Definition~\ref{def:mult}
makes sense over any commutative ring, and the descent works in that generality. For a commutative ring \(R\) of characteristic two,
\(\Ft\subseteq R\), and for an additive map \(\varphi\colon R\to\Ft\) and
\(P\in R[x_\sigma]\) let \(P^\varphi\in\Ft[x_\sigma]\) be obtained by applying
\(\varphi\) to every coefficient.

\begin{lemma}[Descent]\label{lem:descent}
Let \(R\) be a commutative ring of characteristic two, \(\varphi\colon R\to\Ft\)
additive, \(P\in R[x_\sigma]\) and \(a\in\Ft^\sigma\). Every coefficient of
\(P^\varphi(a+z)\) is the image under \(\varphi\) of the corresponding coefficient
of \(P(a+z)\), and \(\deg P^\varphi\le\deg P\)
\lean{claims/PerOrderCharTwo.lean}.
\end{lemma}

\begin{proof}
Write \(P=\sum_uc_ux^u\). Since \(a\) has entries in \(\Ft\), each \((a+z)^u\) has
coefficients \(t_{u,\gamma}\in\Ft\), and the coefficient of \(z^\gamma\) in
\(P(a+z)\) is \(\sum_uc_ut_{u,\gamma}\). For \(t\in\{0,1\}\),
\(\varphi(ct)=t\,\varphi(c)\), so \(\varphi\) maps this coefficient to
\(\sum_u\varphi(c_u)t_{u,\gamma}\), the coefficient of \(z^\gamma\) in
\(P^\varphi(a+z)\). A monomial with coefficient zero in \(P\) has coefficient
\(\varphi(0)=0\) in \(P^\varphi\), so the degree does not increase.
\end{proof}

\begin{theorem}[Every field of characteristic two]\label{thm:char-two}
Let \(R\) be a field of characteristic two, or more generally a commutative ring
of characteristic two, \(n\ge1\) and \(0\le\ell<k\). The least degree of
\(P\in R[x_1,\ldots,x_n]\) with \(\mult_a(P)\ge k\) for every
\(a\in\{0,1\}^n\setminus\{\zero\}\) and \(\mult_{\zero}(P)=\ell\) is
\(\Phi(n,k,\ell)\). For \(k\ge1\), the least degree of such a \(P\) with
\(\mult_{\zero}(P)\le k-1\) is the value \(D(n,k)\) of Corollary~\ref{cor:main}
\lean{claims/PerOrderCharTwo.lean}.
\end{theorem}

\begin{proof}
The polynomial of Theorem~\ref{thm:construction}, read in \(R[x]\), has the same
Hasse coefficients at the points of \(\{0,1\}^n\) and the same degree, because
\(\Ft\to R\) is injective; so the value \(\Phi\) is attained. Conversely, let \(P\)
have multiplicity at least \(k\) at the nonzero points and origin order \(\ell\).
Then \(P\ne0\), and some coefficient \(c\) of \(P(z)\) of degree \(\ell\) is
nonzero. As an \(\Ft\)-vector space \(R\) has an additive map \(\varphi\colon R\to\Ft\)
with \(\varphi(c)\ne0\). By Lemma~\ref{lem:descent}, \(P^\varphi\) has multiplicity
at least \(k\) at the nonzero points, its coefficients of degree below \(\ell\) at
the origin vanish, and its coefficient at the monomial of \(c\) is
\(\varphi(c)\ne0\); so \(\mult_{\zero}(P^\varphi)=\ell\). By Theorem~\ref{thm:main},
\(\deg P\ge\deg P^\varphi\ge\Phi(n,k,\ell)\). The statement about \(D(n,k)\) follows
by minimizing over \(\ell<k\) with Lemma~\ref{lem:minimum}.
\end{proof}

\begin{remark}[Other fields and grids]\label{rem:open}
The proofs of Sections~\ref{sec:upper} and~\ref{sec:lower} use characteristic two
throughout: the expansion in \(x^\varepsilon(x^2+x)^e\) with squarefree
\(x^\varepsilon\), the additivity of \(Y\mapsto Y^{2^j}\), the two shifts
\(\gamma_i\in\{0,1\}\) in Lemma~\ref{lem:congruence}, and the telescoping identity.
The natural next questions are the least degree at every origin order over
fields of odd characteristic and characteristic zero below Menezes' range
\(n\ge k-1\), and the analogue on grids other than \(\{0,1\}^n\).
\end{remark}
